\begin{pastexam}[label={px:ccu-109-swa-06}]{中正資工所甲組 109 學年度 第 11 題}
（追蹤）

(4\%) Please write down the final result of failure array obtained by the following Program (i.e.,
fail function) for each of the following patterns (i.e., pat).
\begin{enumerate}[label=(\alph*)]
\item a b a a b a a b
\item a b c a b a b c a b c
\end{enumerate}
\begin{minted}[fontsize=\small, style=perldoc]{c}
void fail (char *pat, int *failure)
{
  int i, j, n = strlen(pat);
  failure[0] = -1;
  for(j = 1; j<n; j++)
  {
    i = failure[j-1];
    while (pat[j] != pat[i+1] && (i >= 0))
      i = failure[i];
    failure[j] = (pat[j] == pat[i+1]) ? (i+1) : -1;
  }
}
\end{minted}
\end{pastexam}

\begin{pastsoln}
答案：
\begin{center}
\begin{tabular}{l*{11}{c}}
\toprule
$j$ & 0 & 1 & 2 & 3 & 4 & 5 & 6 & 7 & 8 & 9 & 10 \\
\midrule
(a) \texttt{pat[j]} & a & b & a & a & b & a & a & b \\
(a) \texttt{failure[j]} & $-1$ & $-1$ & $0$ & $0$ & $1$ & $2$ & $3$ & $4$ \\
\midrule
(b) \texttt{pat[j]} & a & b & c & a & b & a & b & c & a & b & c \\
(b) \texttt{failure[j]} & $-1$ & $-1$ & $-1$ & $0$ & $1$ & $0$ & $1$ & $2$ & $3$ & $4$ & $2$ \\
\bottomrule
\end{tabular}
\end{center}

\textbf{這個程式的失敗函數是什麼}。 索引從 $0$ 開始，\texttt{failure[j]} 是
\texttt{pat[0..j]} 最長的真邊界（既是真前綴又是後綴的字串）的\textbf{最後一個位置}，也就是邊界長度
減 $1$；沒有邊界時是 $-1$。 它等於資結書第~27.3~節的 $\pi[j+1] - 1$。 迴圈的想法和資結書
虛擬碼~27.3 一樣：先試著把上一格的邊界 \texttt{pat[0..i]} 延長一個字元（比 \texttt{pat[j]} 和
\texttt{pat[i+1]}）；接不上就退到次長的邊界 \texttt{i = failure[i]} 再試，退到 $i = -1$ 為止。
\texttt{while} 條件先比字元再看 $i \ge 0$，$i = -1$ 時比的是 \texttt{pat[0]}，不會越界。

\textbf{(a)} \texttt{abaabaab}，每一格的計算（$i$ 欄是退的過程）：
\begin{center}
\small
\begin{tabular}{cccl}
\toprule
$j$ & \texttt{pat[j]} & $i$ & 結果 \\
\midrule
1 & b & $-1$ & \texttt{pat[0]} $=$ a $\ne$ b，$-1$ \\
2 & a & $-1$ & \texttt{pat[0]} $=$ a，$0$ \\
3 & a & $0 \to -1$ & \texttt{pat[1]} $=$ b $\ne$ a，退到 $-1$；\texttt{pat[0]} $=$ a，$0$ \\
4 & b & $0$ & \texttt{pat[1]} $=$ b，$1$ \\
5 & a & $1$ & \texttt{pat[2]} $=$ a，$2$ \\
6 & a & $2$ & \texttt{pat[3]} $=$ a，$3$ \\
7 & b & $3$ & \texttt{pat[4]} $=$ b，$4$ \\
\bottomrule
\end{tabular}
\end{center}
例如 $j = 7$：\texttt{pat[0..7]} $=$ \texttt{abaabaab} 最長的真邊界是 \texttt{abaab}（長度 $5$），
\texttt{failure[7]} $= 4$。

\textbf{(b)} \texttt{abcababcabc}：
\begin{center}
\small
\begin{tabular}{cccl}
\toprule
$j$ & \texttt{pat[j]} & $i$ & 結果 \\
\midrule
1 & b & $-1$ & \texttt{pat[0]} $=$ a $\ne$ b，$-1$ \\
2 & c & $-1$ & \texttt{pat[0]} $=$ a $\ne$ c，$-1$ \\
3 & a & $-1$ & \texttt{pat[0]} $=$ a，$0$ \\
4 & b & $0$ & \texttt{pat[1]} $=$ b，$1$ \\
5 & a & $1 \to -1$ & \texttt{pat[2]} $=$ c $\ne$ a，退到 \texttt{failure[1]} $= -1$；\texttt{pat[0]} $=$ a，$0$ \\
6 & b & $0$ & \texttt{pat[1]} $=$ b，$1$ \\
7 & c & $1$ & \texttt{pat[2]} $=$ c，$2$ \\
8 & a & $2$ & \texttt{pat[3]} $=$ a，$3$ \\
9 & b & $3$ & \texttt{pat[4]} $=$ b，$4$ \\
10 & c & $4 \to 1$ & \texttt{pat[5]} $=$ a $\ne$ c，退到 \texttt{failure[4]} $= 1$；\texttt{pat[2]} $=$ c，$2$ \\
\bottomrule
\end{tabular}
\end{center}
$j = 10$ 是唯一退了之後又接上的地方：\texttt{pat[0..9]} $=$ \texttt{abcababcab} 的最長邊界
\texttt{abcab} 接不上 c，退到它的最長邊界 \texttt{ab}，\texttt{ab} 後面是 c，接上得 \texttt{abc}，
長度 $3$，\texttt{failure[10]} $= 2$。
\end{pastsoln}
